% Part: incompleteness
% Chapter: arithmetization-syntax
% Section: proofs-in-ax

\documentclass[../../../include/open-logic-section]{subfiles}

\begin{document}

\olfileid{inc}{art}{pax}
\olsection{स्वयंसिद्धकीय \usetoken{P}{derivation}}

\begin{explain}
स्वयंसिद्धकीय !!{derivation}s चे अंकगणितीकरण करण्यासाठी
!!{derivation}s संख्यांनी दर्शवावे लागतात. !!{derivation}s या
!!{formula}s च्या क्रमिका असल्यामुळे, प्रत्येक !!{derivation}
तिच्यातील !!{formula}s च्या संकेतांकांच्या क्रमिकेचा संकेतांक
म्हणून दर्शवणे हा सहज मार्ग आहे.
\end{explain}

\begin{defn}
$\delta$ ही $!A_1$, \dots,~$!A_n$ या !!{formula}s ची
स्वयंसिद्धकीय !!{derivation} असेल, तर $\Gn{\delta}$ म्हणजे
\[
\tuple{\Gn{!A_1}, \dots, \Gn{!A_n}}.
\]
\end{defn}

\begin{ex}
  पुढील अगदी सोपी !!{derivation} पाहा:
  \begin{derivation}
  1. & $!B \lif (!B \lor !A)$ \\
  2. & $(!B \lif (!B \lor !A)) \lif (!A  \lif (!B \lif (!B \lor !A)))$\\
  3. & $!A  \lif (!B \lif (!B \lor !A))$
  \end{derivation}
  या !!{derivation} ची गोडेल संख्या
  \begin{align*}
  \openTuple\,
    & \Gn{!B \lif (!B \lor !A)}, \\
    &\Gn{(!B \lif (!B \lor !A)) \lif (!A  \lif (!B \lif (!B \lor !A)))},\\
    & \Gn{!A  \lif (!B \lif (!B \lor !A))} \,\closeTuple.
  \end{align*}
  अशी असेल.
\end{ex}

\begin{explain}
!!{derivation}s चे निरूपण ठरल्यानंतर, अशा !!{derivation}s वर
आदिम पुनरावर्ती रीतीने क्रिया करता येतात आणि त्यांचे आवश्यक
गुणधर्म व संबंध तसेच व्यक्त करता येतात, हेही दाखवावे लागेल.
काही क्रिया सोप्या आहेत. उदाहरणार्थ, !!a{derivation} ची गोडेल
संख्या~$d$ दिली असता, $(d)_{\len{d}-1}$ हे तिच्या अंतिम
!!{formula} ची गोडेल संख्या देते. काही इतर क्रिया अधिक कठीण
आहेत; त्या कशा कराव्यात याची किमान रूपरेषा आपण पाहू.
``$\delta$ ही $\Gamma$ पासून $!A$ ची !!a{derivation} आहे''
हा संबंध $\delta$ आणि $!A$ यांच्या गोडेल संख्यांचा आदिम
पुनरावर्ती संबंध आहे, हे दाखवणे आपले उद्दिष्ट आहे.
\end{explain}

\begin{prop}
\ollabel{prop:followsby}
पुढील संबंध आदिम पुनरावर्ती आहेत:
\begin{enumerate}
\item $!A$ हे स्वयंसिद्धक आहे.
\item $\delta$ मधील $i$ वी ओळ विध्यात्मक अनुमानाने समर्थित आहे.
\item $\delta$ मधील $i$ वी ओळ \QR{} ने समर्थित आहे.
\item $\delta$ ही योग्य !!{derivation} आहे.
\end{enumerate}
\end{prop}

\begin{proof}
!!{formula}s च्या गोडेल संख्या आणि !!{derivation}s च्या
गोडेल संख्या यांच्यातील अनुरूप संबंध आदिम पुनरावर्ती आहेत,
हे दाखवायचे आहे.
\begin{enumerate}
\item स्वयंसिद्धक-रूपबंधांची यादी आपल्याला दिली आहे; $!A$ हे
  स्वयंसिद्धक असेल, तर ते या रूपबंधांपैकी एखाद्याच्या रूपाचे असते.
  यादी सांत असल्यामुळे, प्रत्येक रूपबंधासाठी $!A$ त्या रूपाचे
  आहे का हे आदिम पुनरावर्ती रीतीने तपासता येते, एवढे दाखवणे पुरेसे आहे.
  उदाहरणार्थ पुढील स्वयंसिद्धक-रूपबंध पाहा:
  \[
  !B \lif (!C \lif !B).
  \]
  $!B$ आणि $!C$ ही !!{formula}s अशी असतील की `$($', $!B$,
  `$\lif$', `$($', $!C$, `$\lif$', $!B$ आणि~`$))$' या
  चिन्हमाला जोडून $!A$ मिळते, तर $!A$ हा या रूपबंधाचा नमुना आहे.
  $!A$ ची गोडेल संख्या $n$ असल्यास हा गुणधर्म तपासता येतो:
  क्रमिकांची जोडणी आदिम पुनरावर्ती आहे आणि हा संबंध खरा असेल,
  तर $!B$ व $!C$ ही $!A$ ची उप-!!{formula}s असल्यामुळे $!B$ व $!C$
  यांच्या गोडेल संख्या $!A$ च्या गोडेल संख्येपेक्षा लहान असतात.
  म्हणून व्याख्या अशी करता येते:
  \begin{multline*}
  \fn{IsAx}_{!B \lif (!C \lif !B)}(n) \defiff \bexists{b< n}{
    \bexists{c < n}{(\fn{Frm}(b) \land \fn{Frm}(c) \land {}}}\\ n =
  \Gn{(} \concat b \concat \Gn{\lif} \concat \Gn{(} \concat c \concat
  \Gn{\lif} \concat b \concat \Gn{))}).
  \end{multline*}
  प्रत्येक स्वयंसिद्धक-रूपबंधासाठी अशी व्याख्या केली, तर
  त्यांच्या विकल्पाने $\fn{IsAx}(n)$ हा गुणधर्म परिभाषित होतो:
  ``$n$ ही स्वयंसिद्धकाची गोडेल संख्या आहे.''
\item $\delta$ मधील $i$ वी ओळ विध्यात्मक अनुमानाने समर्थित असेल,
  तर आणि तरच $j$ आणि $k < i$ अशा आधीच्या ओळी असतात की
  $j$ व्या ओळीवरील !!{formula}~$!A$ असते,
  $k$ व्या ओळीवरील सूत्र $!A \lif !B$ असते, आणि
  $i$ व्या ओळीवरील सूत्र~$!B$ असते.
  \begin{multline*}
    \fn{MP}(d, i) \defiff \bexists{j < i}{\bexists{k < i}{}}\\
    (d)_k = \Gn{(} \concat (d)_j
      \concat \Gn{\lif} \concat (d)_i \concat \Gn{)}
  \end{multline*}
  परिबद्ध संख्यकीकरण, जोडणी आणि $=$ हे आदिम पुनरावर्ती
  असल्यामुळे हा संबंधही आदिम पुनरावर्ती आहे.
\item $\delta$ मधील एखादी ओळ \QR{} ने समर्थित असेल,
  तर ती $!B \lif \lforall[x][!A(x)]$ या रूपाची असते;
  आधीची एखादी ओळ, कोणत्यातरी !!{constant}~$c$ साठी,
  $!B \lif !A(c)$ असते; आणि $c$ हा $!B$ मध्ये येत नाही.
  पुढील अटी पूर्ण झाल्या तर आणि तरच असे घडते:
  \begin{enumerate}
  \item एखादे !!a{sentence}~$!B$ असते (येथे हा विशेष प्रसंग;
    सर्वसाधारण नियमात सूत्रही चालते), आणि
  \item $x$ हाच मुक्त चल असू शकणारे !!a{formula}~$!A(x)$ असते, ज्यासाठी
  \item ओळ $i$ मध्ये $!B \lif \lforall[x][!A(x)]$ असते,
  \item आधीच्या एखाद्या ओळीत, $j < i$, कोणत्यातरी अचर~$c$ साठी
    $!B \lif \Subst{!A}{c}{x}$ असते,
  \item आणि हा अचर $!B$ मध्ये येत नाही.
  \end{enumerate}
  या अटी आदिम पुनरावर्ती रीतीने तपासता येतात. कारण $!B$,
  $!A(x)$ आणि $x$ यांच्या गोडेल संख्या $i$ व्या ओळीवरील
  सूत्राच्या गोडेल संख्येपेक्षा लहान आहेत; तसेच $c$ ची गोडेल
  संख्या $j$ व्या ओळीवरील सूत्राच्या गोडेल संख्येपेक्षा लहान आहे:
  \begin{multline*}
    \fn{QR}_1(d, i) \defiff \bexists{j<i}{\bexists{b < (d)_i}{\bexists{x <
        (d)_i}{\bexists{a < (d)_i}{\bexists{c < (d)_j}{(}}}}}
    \\ \fn{Var}(x) \land \fn{Const}(c) \land {} \\
    (d)_i = \Gn{(} \concat
    b \concat \Gn{\lif} \concat \Gn{\lforall} \concat x \concat a
    \concat \Gn{)} \land {}\\
    (d)_j = \Gn{(} \concat b \concat
    \Gn{\lif} \concat \fn{Subst}(a,c,x) \concat \Gn{)} \land {}\\ \fn{Frm}(b)
    \land \fn{Frm}(a) \land {} \bforall{k <
      \len{b}}{(b)_k \neq (c)_0})
  \end{multline*}
  येथे $c$ आणि $x$ यांना अचर आणि चल यांच्या पदरूपातील
  गोडेल संख्या मानले आहे (म्हणजे त्यांचे चिन्हसंकेतांक नाहीत).
  $x$ हा $!A(x)$ मधील आदेशनासाठीचा चल आहे;
  $\Subst{!A(x)}{c}{x}$ हे !!a{sentence} आहे का ही मूळ
  वाक्यांपुरती तपासणी सर्वसाधारण संख्यक-नियमासाठी आवश्यक नाही.
  सूत्ररूपता तपासून, $!B$ मधील प्रत्येक चिन्ह $c$ पेक्षा
  वेगळे आहे ही अट ठेवून $c$ हा $!B$ मध्ये येत नाही,
  याची खात्री करतो.

  \QR{} ची दुसरी आवृत्ती सरावासाठी ठेवतो.
\item $d$ ही योग्य !!{derivation} ची गोडेल संख्या असेल, तर आणि
  तरच ती रिकामी नसते आणि तिच्यातील प्रत्येक ओळ स्वयंसिद्धक
  किंवा विध्यात्मक अनुमान अथवा \QR{} ने समर्थित असते. म्हणून:
  \[
  \fn{Deriv}(d) \defiff \len{d}>0 \land
  \bforall{i < \len{d}}{(\fn{IsAx}((d)_i) \lor
    \fn{MP}(d,i) \lor \fn{QR}(d, i))}
  \]
\end{enumerate}
\end{proof}

\begin{prob}
\olref[inc][art][pax]{prop:followsby} प्रमाणे पुढील संबंध
परिभाषित करा:
\begin{enumerate}
\item $\fn{IsAx}_{!A \lif (!B \lif (!A \land !B))}(n)$,
\item $\fn{IsAx}_{\lforall[x][!A(x)] \lif !A(t)}(n)$,
\item $\fn{QR}_{2}(d, i)$ (\QR{} च्या दुसऱ्या आवृत्तीसाठी).
\end{enumerate}
\end{prob}

\begin{prop}
समजा $\Gamma$ हा !!{sentence}s चा आदिम पुनरावर्ती संच आहे.
मग $\Prf[\Gamma](x, y)$ हा संबंध आदिम पुनरावर्ती आहे.
तो असे व्यक्त करतो: $x$ हा $\Gamma$ पासून $!A$ च्या
!!a{derivation}~$\delta$ चा संकेतांक आहे आणि $y$ ही
$!A$ ची गोडेल संख्या आहे.
\end{prop}

\begin{proof}
समजा ``$y \in \Gamma$'' हे $R_\Gamma(y)$ या आदिम पुनरावर्ती
विधेयाने दिले आहे. $\Prf[\Gamma](x, y)$ हा संबंध आदिम पुनरावर्ती
आहे, हे दाखवायचे आहे. $y$ ही !!a{sentence}~$!A$ ची गोडेल संख्या
असेल आणि $x$ हा $\Gamma$ पासून $!A$ च्या
!!a{derivation} चा संकेतांक असेल, तर आणि तरच
$\Prf[\Gamma](x, y)$ खरा असतो.

मागील प्रस्तावानुसार, $x$ हा योग्य !!{derivation}~$\delta$
चा संकेतांक असेल, तर आणि तरच खरा असणारा $\fn{Deriv}(x)$
गुणधर्म आदिम पुनरावर्ती आहे. मात्र त्या व्याख्येत
!!{derivation} मधील ओळी समर्थित करण्याचा अतिरिक्त मार्ग म्हणून
$\Gamma$ संच विचारात घेतलेला नाही. ओळ \QR{} ने समर्थित आहे
का हे तपासणाऱ्या आदिम पुनरावर्ती कसोटीतही अचर~$c$ हा
$\Gamma$ मध्ये येऊ नये ही अट घेतलेली नाही. $\Gamma$ चा
थेट विचार करून व्याख्या सुधारता येईल; पण $\fn{Deriv}$ आणि
निगमन प्रमेय वापरणे सोपे आहे. काही सांत !!{sentence}s ची यादी
$!B_1$, \dots, $!B_n \in \Gamma$ अशी असेल की
$\{!B_1, \dots, !B_n\} \Proves !A$, तर आणि तरच
$\Gamma \Proves !A$. निगमन प्रमेयानुसार, हे
$\Proves (!B_1 \lif (!B_2 \lif
\cdots (!B_n \lif !A)\cdots))$ असेल तेव्हा घडते.
गोडेल संख्या~$z$ असलेले !!a{sentence} या रूपाचे आहे का,
हे आदिम पुनरावर्ती रीतीने तपासता येते. म्हणून $x$ ला
गोडेल संख्या~$y$ असलेल्या !!{sentence} ची \emph{$\Gamma$ पासून}
!!a{derivation} असण्याचा थेट संकेतांक मानण्याऐवजी,
$x$ हा वरच्या रूपातील आवर्ती सोपाधिक वाक्याच्या
$\emptyset$ पासूनच्या !!a{derivation} चा संकेतांक मानतो.
यामुळे खालील कसोटी सिद्धतायोग्यतेसाठी समतुल्य असली,
तरी संकेतांकाचा अर्थ सुरुवातीच्या विधानातील थेट संकेतांकापेक्षा बदलतो.

प्रथम, !!{sentence}s ची क्रमिका दिल्यास तिच्यातील सर्व वाक्ये
पूर्वांग म्हणून आणि दिलेले !!{sentence} उत्तरांग म्हणून असलेले
सोपाधिक वाक्य आदिम पुनरावर्ती रीतीने तयार करता येते:
\begin{align*}
  \fn{hCond}(s, y, 0) & = y \\
  \fn{hCond}(s, y, n+1) & = \Gn{(} \concat (s)_{n} \concat \Gn{\lif}
  \concat \fn{hCond}(s, y, n) \concat \Gn{)}\\
  \fn{Cond}(s, y) & = \fn{hCond}(s, y, \len{s})\\
  \intertext{म्हणून $\Prf[\Gamma](x, y)$ ची व्याख्या अशी करता येते:}
  \Prf[\Gamma](x, y) & \defiff \fn{Sent}(y) \land
  \bexists{s < \fn{sequenceBound}(x,x)}{(} \\
  &\qquad (x)_{\len{x}-1} = \fn{Cond}(s,y) \land {} \\
  &\qquad\bforall{i<\len{s}}{(s)_i \in \Gamma} \land {} \\
  &\qquad\fn{Deriv}(x)).
\end{align*}
येथे पुनरावर्तनाने क्रमिकेतील पूर्वांगे आतून बाहेर या क्रमाने
जोडली जातात; आधी लिहिलेल्या बाहेरून आतच्या यादीसाठी ती उलट क्रमाने
निवडता येतात. $s$ वरील परिबंध असा मिळतो: प्रत्येक $(s)_i$ ही
!!{derivation} च्या अखेरच्या ओळीतील उप-!!{formula} ची गोडेल संख्या
आहे; म्हणजेच ती $(x)_{\len{x}-1}$ पेक्षा लहान आहे. $!B
\in \Gamma$ या पूर्वांगांची संख्या, म्हणजे $s$ ची लांबी,
$x$ च्या अखेरच्या ओळीच्या लांबीपेक्षा कमी आहे.
\end{proof}

\end{document}
